\documentclass[10pt,a4paper]{amsart}
\usepackage{amsmath,amssymb,amsthm,graphicx}
\usepackage{verbatim}
\newenvironment{reference}{\comment}{\endcomment}
%\newenvironment{reference}{}{}
\newenvironment{slogan}{\comment}{\endcomment}
\newenvironment{history}{\comment}{\endcomment}

% For commutative diagrams we use Xy-pic
\usepackage[all]{xy}

% We use 2cell for 2-commutative diagrams.
\xyoption{2cell}
% Theorem environments.
%
\theoremstyle{plain}
\newtheorem{theorem}[subsection]{Theorem}
\newtheorem{proposition}[subsection]{Proposition}
\newtheorem{lemma}[subsection]{Lemma}

\theoremstyle{definition}
\newtheorem{definition}[subsection]{Definition}
\newtheorem{example}[subsection]{Example}
\newtheorem{exercise}[subsection]{Exercise}
\newtheorem{situation}[subsection]{Situation}

\theoremstyle{remark}
\newtheorem{remark}[subsection]{Remark}
\newtheorem{remarks}[subsection]{Remarks}

\numberwithin{equation}{subsection}

% Macros
%
\def\lim{\mathop{\mathrm{lim}}\nolimits}
\def\colim{\mathop{\mathrm{colim}}\nolimits}
\def\Spec{\mathop{\mathrm{Spec}}}
\def\Hom{\mathop{\mathrm{Hom}}\nolimits}
\def\Ext{\mathop{\mathrm{Ext}}\nolimits}
\def\SheafHom{\mathop{\mathcal{H}\!\mathit{om}}\nolimits}
\def\SheafExt{\mathop{\mathcal{E}\!\mathit{xt}}\nolimits}
\def\Sch{\mathit{Sch}}
\def\Mor{\mathop{\mathrm{Mor}}\nolimits}
\def\Ob{\mathop{\mathrm{Ob}}\nolimits}
\def\Sh{\mathop{\mathit{Sh}}\nolimits}
\def\NL{\mathop{N\!L}\nolimits}
\def\CH{\mathop{\mathrm{CH}}\nolimits}
\def\proetale{{pro\text{-}\acute{e}tale}}
\def\etale{{\acute{e}tale}}
\def\QCoh{\mathit{QCoh}}
\def\Ker{\mathop{\mathrm{Ker}}}
\def\Im{\mathop{\mathrm{Im}}}
\def\Coker{\mathop{\mathrm{Coker}}}
\def\Coim{\mathop{\mathrm{Coim}}}

% Boxtimes
%
\DeclareMathSymbol{\boxtimes}{\mathbin}{AMSa}{"02}

%
% Macros for moduli stacks/spaces
%
\def\QCohstack{\mathcal{QC}\!\mathit{oh}}
\def\Cohstack{\mathcal{C}\!\mathit{oh}}
\def\Spacesstack{\mathcal{S}\!\mathit{paces}}
\def\Quotfunctor{\mathrm{Quot}}
\def\Hilbfunctor{\mathrm{Hilb}}
\def\Curvesstack{\mathcal{C}\!\mathit{urves}}
\def\Polarizedstack{\mathcal{P}\!\mathit{olarized}}
\def\Complexesstack{\mathcal{C}\!\mathit{omplexes}}
% \Pic is the operator that assigns to X its picard group, usage \Pic(X)
% \Picardstack_{X/B} denotes the Picard stack of X over B
% \Picardfunctor_{X/B} denotes the Picard functor of X over B
\def\Pic{\mathop{\mathrm{Pic}}\nolimits}
\def\Picardstack{\mathcal{P}\!\mathit{ic}}
\def\Picardfunctor{\mathrm{Pic}}
\def\Deformationcategory{\mathcal{D}\!\mathit{ef}}
\usepackage{cleveref}
\crefname{originalsection}{original section}{original sections}
\Crefname{originalsection}{Original section}{Original sections}
\makeatletter
\newlabel{sites-section-presheaves}{{00V1}{source}{}{sourcepointer-1}{}}
\newlabel{sites-section-presheaves@cref}{{[originalsection][0][]00V1}{[1][0][]source}{}{}{}}
\makeatother
\crefname{originalsection}{original section}{original sections}
\Crefname{originalsection}{Original section}{Original sections}
\makeatletter
\newlabel{sites-section-representable-sheaves}{{00WO}{source}{}{sourcepointer-2}{}}
\newlabel{sites-section-representable-sheaves@cref}{{[originalsection][0][]00WO}{[1][0][]source}{}{}{}}
\makeatother
\crefname{originalproposition}{original proposition}{original propositions}
\Crefname{originalproposition}{Original proposition}{Original propositions}
\makeatletter
\newlabel{proposition-etale-morphisms}{{03PC}{source}{}{sourcepointer-3}{}}
\newlabel{proposition-etale-morphisms@cref}{{[originalproposition][0][]03PC}{[1][0][]source}{}{}{}}
\makeatother
\title{Global sections of a sheaf on the etale site of a field are the Galois invariants of its stalk}
\author{The Stacks Project Authors}
\date{stacks-project:etale-cohomology.tex:local-snapshot-20261011; GNU FDL 1.2 or later}
\begin{document}
\maketitle
\noindent\textit{Editorial note.} All delivered mathematical text is the author's own, copied from the pinned Stacks Project chapter file etale-cohomology.tex; no statement or proof step was rewritten, reordered or removed.\par\medskip
\noindent\textit{Editorial note.} The proof is the author's complete argument and even offers a second, direct proof; nothing but the identification of the last step of the formal chain as clear is left to the reader.\par\medskip
\noindent\textit{Editorial note.} Every labelled reference occurring inside the counted spans is carried in assembly.external\_references as a non-load-bearing navigation pointer (same-chapter ones at their definition site, cross-chapter ones at the invocation inside the counted span), so no load-bearing mathematics is replaced by a stub. No bibliography citation occurs inside the counted spans.\par\medskip
\noindent\textit{Editorial note.} The checked chapter list for same-name collisions is the set of stacks chapters staged in this repository: chow, cohomology, divisors, etale-cohomology, flat, more-algebra, more-morphisms, morphisms, spaces-more-morphisms, varieties together with this card's chapter etale-cohomology.\par\medskip
\noindent\textit{Editorial note.} The preamble is the chapter's own native preamble excerpt already used for this source: the Xy-pic setup (native preamble lines 14-18) and the native theorem declarations with the mathematical macros (native preamble lines 154-217).\par\medskip
\noindent\textit{Allowed standard prerequisites (editorial).}\begin{itemize}
\item Equivalence between sheaves of sets and Galois sets. Same chapter, Theorem theorem-equivalence-sheaves-point. Turns the given sheaf into a G-set whose invariants compute the global sections.
\item Sheaves represented by objects of the site. Sites, Sections sites-section-presheaves and sites-section-representable-sheaves. Identifies the global sections with morphisms from the sheaf represented by the site's final object.
\end{itemize}\medskip
\ifdefined\hypertarget\hypertarget{sourcepointer-1}{}\fi
\noindent\textit{Source reference.} Original section 00V1: sources/stacks-etale-cohomology/pinned-original-source.tex, source lines 8690--8690. Navigation pointer invoked at proof line 8690; the counted span contains the pointer only and the cited result is not reproduced inside it, so the pointer is registered as a non-load-bearing navigation stub in the established form for Stacks material. No bibliography citation is involved.\par\medskip
\ifdefined\hypertarget\hypertarget{sourcepointer-2}{}\fi
\noindent\textit{Source reference.} Original section 00WO: sources/stacks-etale-cohomology/pinned-original-source.tex, source lines 8691--8691. Navigation pointer invoked at proof line 8691; the counted span contains the pointer only and the cited result is not reproduced inside it, so the pointer is registered as a non-load-bearing navigation stub in the established form for Stacks material. No bibliography citation is involved.\par\medskip
\ifdefined\hypertarget\hypertarget{sourcepointer-3}{}\fi
\noindent\textit{Source reference.} Original proposition 03PC: sources/stacks-etale-cohomology/pinned-original-source.tex, source lines 2592--2592. Navigation pointer invoked at proof line 8705; the counted span contains the pointer only and the cited result is not reproduced inside it, so the pointer is registered as a non-load-bearing navigation stub in the established form for Stacks material. No bibliography citation is involved.\par\medskip
\begin{theorem}
\label{theorem-equivalence-sheaves-point}
Let $S = \Spec(K)$ with $K$ a field.
Let $\overline{s}$ be a geometric point of $S$.
Let $G = \text{Gal}_{\kappa(s)}$ denote the absolute Galois group.
Taking stalks induces an equivalence of categories
$$
\Sh(S_\etale) \longrightarrow G\textit{-Sets},
\quad
\mathcal{F} \longmapsto \mathcal{F}_{\overline{s}}.
$$
\end{theorem}

\begin{lemma}
\label{lemma-global-sections-point}
Assumptions and notations as in
Theorem \ref{theorem-equivalence-sheaves-point}.
There is a functorial bijection
$$
\Gamma(S, \mathcal{F}) = (\mathcal{F}_{\overline{s}})^G
$$
\end{lemma}

\begin{proof}
We can prove this using formal arguments and the result of
Theorem \ref{theorem-equivalence-sheaves-point}
as follows. Given a sheaf $\mathcal{F}$ corresponding to
the $G$-set $M = \mathcal{F}_{\overline{s}}$ we have
\begin{eqnarray*}
\Gamma(S, \mathcal{F}) & = &
\Mor_{\Sh(S_\etale)}(h_{\Spec(K)},  \mathcal{F})
\\
& = & \Mor_{G\textit{-Sets}}(\{*\}, M) \\
& = & M^G
\end{eqnarray*}
Here the first identification is explained in
Sites, Sections \ref{sites-section-presheaves} and
\ref{sites-section-representable-sheaves},
the second results from
Theorem \ref{theorem-equivalence-sheaves-point}
and the third is clear. We will also give a direct proof\footnote{For
the doubting Thomases out there.}.

\medskip\noindent
Suppose that $t \in \Gamma(S, \mathcal{F})$ is a global section.
Then the triple $(S, \overline{s}, t)$ defines an element of
$\mathcal{F}_{\overline{s}}$ which is clearly invariant under the
action of $G$. Conversely, suppose that $(U, \overline{u}, t)$
defines an element of $\mathcal{F}_{\overline{s}}$ which is invariant.
Then we may shrink $U$ and assume $U = \Spec(L)$ for some
finite separable field extension of $K$, see
Proposition \ref{proposition-etale-morphisms}.
In this case the map $\mathcal{F}(U) \to \mathcal{F}_{\overline{s}}$
is injective, because for any morphism of \'etale neighbourhoods
$(U', \overline{u}') \to (U, \overline{u})$ the restriction map
$\mathcal{F}(U) \to \mathcal{F}(U')$ is injective since $U' \to U$
is a covering of $S_\etale$.
After enlarging $L$ a bit we may assume $K \subset L$ is a finite
Galois extension. At this point we use that
$$
\Spec(L) \times_{\Spec(K)} \Spec(L)
=
\coprod\nolimits_{\sigma \in \text{Gal}(L/K)} \Spec(L)
$$
where the maps $\Spec(L) \to \Spec(L \otimes_K L)$
come from the ring maps $a \otimes b \mapsto a\sigma(b)$. Hence we
see that the condition that $(U, \overline{u}, t)$ is invariant
under all of $G$ implies that $t \in \mathcal{F}(\Spec(L))$
maps to the same element of
$\mathcal{F}(\Spec(L) \times_{\Spec(K)} \Spec(L))$
via restriction by either projection (this uses the injectivity mentioned
above; details omitted). Hence the sheaf condition of $\mathcal{F}$
for the \'etale covering $\{\Spec(L) \to \Spec(K)\}$ kicks
in and we conclude that $t$ comes from a unique section of $\mathcal{F}$
over $\Spec(K)$.
\end{proof}

\end{document}
